We study sparse traffic sensor siting as a recoverability-driven network design problem. Given a fixed sensor budget, the goal is to choose a long-lived layout that minimizes hidden-state reconstruction error under a transparent generalized-least-squares/maximum-a-posteriori estimator. We formulate TRACE-BiOpt, a bilevel stochastic design method whose lower level reconstructs the full traffic state and whose upper level combines hidden-state Huber reconstruction loss, posterior uncertainty, conditional-value-at-risk tail risk, and spatial redundancy regularization. A deterministic solver uses scale-adaptive initialization and one-exchange refinement under the same objective; pre-registered baselines are not used as method candidates. We prove lower-level stability, a posterior-risk identity, an all-layout validation generalization bound, and an exchange-stationarity certificate. Across PeMS7\_228, PeMS7\_1026, and Seattle at 10\%, 20\%, and 30\% budgets over ten split seeds, TRACE-BiOpt achieves the lowest mean held-out GLS/MAP MAE against 21 pre-registered non-BiOpt baselines on all nine dataset-budget rows. Holm-corrected paired tests leave no statistically tied or better pre-registered challenger. The claims are scoped to the tested regimes and do not assert global layout optimality, universal robustness, or dominance over untested methods.
